% Generated draft: requires mathematical and source review.
\documentclass[11pt]{article}
\usepackage[margin=25mm]{geometry}
\usepackage{fontspec}
\setmainfont{Latin Modern Roman}
\usepackage{amsmath,amssymb,mathtools,unicode-math}
\setmathfont{Latin Modern Math}
\usepackage{enumitem,needspace,xurl,hyperref}
\setlength{\parindent}{0pt}
\setlength{\parskip}{5pt}
\setlength{\emergencystretch}{3em}
\newcommand{\N}{\mathbb{N}}
\newcommand{\Z}{\mathbb{Z}}
\newcommand{\Q}{\mathbb{Q}}
\newcommand{\R}{\mathbb{R}}
\newcommand{\C}{\mathbb{C}}
\newcommand{\F}{\mathbb{F}}
\newcommand{\A}{\mathbb{A}}
\newcommand{\PP}{\mathbb{P}}
\newcommand{\E}{\mathbb{E}}
\newcommand{\eps}{\varepsilon}
\newcommand{\dd}{\,\mathrm{d}}
\begin{document}
\section*{UnsolvedMath ID: 30000460}
\textbf{Source problem code:} catalogue entry
\subsection{Definitions and mathematical statement}
Let \(\mathcal O\) be the coefficient ring intended in the source, let \(\mathfrak S_n\) be a finite symmetric group, and let \(b\) be a primitive central block idempotent of \(\mathcal O\mathfrak S_n\). Write \(B=\mathcal O\mathfrak S_n b\) for the block algebra. The conjectural classification asks that every such \(B\) be equivalent to a quotient of the deformation of the double of the Schur-algebra object associated to \(\mathcal O Q\), and also to the analogous quotient in the Koszul-dual construction.

Both quotient descriptions are required, not an alternative between them. The exact algebra \(\mathcal O Q\), its Schur construction and degree parameters, the double, deformation family and specialization, quotient ideals, and Koszul-dual antecedent are not defined in the record. In particular \(Q\) could be source quiver data; it must not be read automatically as the rational-number field, and the double is not silently identified with a Drinfeld double, trivial extension, or doubled-quiver algebra.

The word equivalent also requires specification. Under a Morita reading, one would require an equivalence of module categories
\[
B\text{-}\operatorname{mod}\simeq(D/I)\text{-}\operatorname{mod}
\]
for the intended deformed-double algebra \(D\) and quotient ideal \(I\), and a second analogous equivalence for the dual construction. An algebra isomorphism, derived equivalence, stable equivalence, or merely equality of numerical invariants would be different claims. This display explains the Morita possibility without choosing it as the unprovided source convention. Any grading and \(\mathcal O\)-linearity conditions must also be recovered.

The source does not give a characteristic or splitting system for \(\mathcal O\), a block-weight restriction, or a rank restriction on the symmetric groups. None is introduced from a guessed familiar construction. The supplied literature assessment checked 2026-08-21 calls the all-blocks deformed-double description a broad conjectural classification and reports that no general proof was found. Its exact constructions must be retrieved before there is a fully defined quotient theorem.
\subsection{Short English statement}
Conjecturally every symmetric-group block has both proposed deformed-double Schur-algebra quotient models, including the dual version. The algebras, deformation and quotient parameters, and meaning of equivalence are missing, so those descriptions cannot yet be made into one precise classification theorem.
\subsection{Sources}
\begin{itemize}
\item Local UnsolvedMath registry, permanent ID 30000460, catalogue entry, \emph{Schur-Algebra Models for Symmetric-Group Blocks}.
\item \emph{Representations of Finite Groups} (2007), Oberwolfach Reports, pp. 895--978, DOI \url{https://doi.org/10.4171/owr/2006/15}.
\item Supplied literature triage checked 2026-08-21 cites only the original report; no separate theorem or complete definition of the proposed Schur, double, deformation, or dual construction is supplied.
\end{itemize}
\end{document}
